\documentclass[12pt]{article}
\usepackage[utf8]{inputenc}
\usepackage[spanish]{babel}
\usepackage{csquotes}
\usepackage{geometry}
\usepackage{graphicx}
\usepackage{hyperref}
\geometry{a4paper, margin=2.5cm}

\title{Geopet: una solución interactiva y didáctica para el turismo geolocalizado}
\author{Antony Quispe Calizaya, Brayan Calderon Calderon}
\date{\today}

\begin{document}
\maketitle

\begin{abstract}
Geopet es una aplicación móvil diseñada para promover el turismo mediante una experiencia interactiva y gamificada. Este artículo presenta el problema actual en la experiencia turística moderna, caracterizada por baja interactividad, poca retención de información cultural, desinformación logística y la ausencia de herramientas tecnológicas en muchas regiones con potencial turístico.

Se exploran las carencias de guías y visitas  que no logran cautivar ni informar adecuadamente al visitante moderno, especialmente a las nuevas generaciones de turistas digitales. En este contexto, Geopet propone una solución innovadora que combina tecnologías geoespaciales (GIS), elementos de gamificación y estrategias de edutainment para transformar la exploración de espacios culturales y naturales en una experiencia significativa y personalizada.

A través de rutas inteligentes, recompensas digitales, narrativa interactiva y visualización jerárquica de puntos de interés, Geopet no solo mejora la navegación y el acceso a la información, sino que también fortalece la retención del conocimiento y el vínculo emocional con el territorio. Este artículo analiza los fundamentos conceptuales del diseño de la aplicación, sus beneficios esperados y las implicancias para el desarrollo sostenible del turismo en regiones subatendidas.
\end{abstract}

\vspace{1cm}
\section{Introducción}
En el contexto del turismo contemporáneo, los visitantes enfrentan múltiples desafíos: saturación de información, itinerarios poco personalizados y escasa motivación para explorar más allá de los puntos emblemáticos. Estudios señalan que los turistas buscan experiencias significativas que combinen entretenimiento, conocimiento y facilidad de uso \cite{Xu2015,Pradhan2023}. Sin embargo, las guías tradicionales y aplicaciones convencionales no siempre satisfacen estas expectativas.

Con la creciente digitalización del turismo, los usuarios demandan experiencias interactivas, educativas y accesibles a través de sus dispositivos móviles \cite{Fishoeder2018,Frontiers2024}. La gamificación, definida como el uso de mecánicas de juego en contextos no lúdicos, ha demostrado ser efectiva para mejorar el compromiso del usuario, aumentar la retención de información y fomentar la exploración activa \cite{Weber2017,Swacha2017}. A pesar de su eficacia, muchas aplicaciones turísticas siguen siendo estáticas, centradas en mapas básicos o en listados de lugares sin personalización ni interacción \cite{Xu2014,Architecture2019}.

En particular, en muchas regiones emergentes o con un patrimonio turístico en desarrollo, como es el caso de la región de estudio propuesta, se observa una notoria ausencia de soluciones tecnológicas. Actualmente, los visitantes dependen exclusivamente de guías turísticos presenciales, sin el apoyo de folletos gubernamentales, plataformas web interactivas o aplicaciones móviles especializadas. Esta carencia limita la autonomía, la planificación y la exploración consciente del entorno \cite{ToARist2020,Emerald2020,Intech2019}.

La implementación de herramientas basadas en Sistemas de Información Geográfica (GIS) puede ser clave para solucionar estos vacíos. Interfaces visuales e interactivas permiten representar datos espaciales de forma intuitiva, facilitando la navegación, la selección de rutas optimizadas y la jerarquización de puntos de interés según las preferencias del usuario \cite{Arxiv2025,Kawanaka2021,NatureGIS2022}. Además, las estrategias de edutainment han sido exitosas en mejorar la retención del conocimiento cultural, al integrar narrativas, recompensas y quizzes dentro de la experiencia \cite{Frontiers2024,TravelPlot2018}.

Desde esta perspectiva, surge la necesidad de desarrollar soluciones que integren gamificación, GIS y diseño centrado en el usuario para ofrecer experiencias turísticas personalizadas, educativas y sostenibles. Este artículo se enfoca en presentar una propuesta que responde a esa necesidad: la aplicación Geopet.

\section{El problema}
El turismo, como actividad multidimensional, requiere herramientas modernas que integren tecnologías geoespaciales, interactividad y personalización. En ausencia de estas herramientas, los turistas se enfrentan a experiencias fragmentadas, con dificultades para obtener información relevante, navegar eficientemente entre puntos de interés, y conectar emocionalmente con los lugares visitados \cite{Pradhan2023,Weber2017}. Esta situación es particularmente crítica en contextos regionales donde no existen mecanismos digitales ni estrategias de gamificación implementadas, lo cual limita el potencial de desarrollo turístico sostenible e inclusivo.

\subsection{Falta de interactividad y motivación activa}
Las aplicaciones turísticas tradicionales, si bien han digitalizado parte de la experiencia, en su mayoría ofrecen contenido estático: listas de puntos de interés, imágenes y descripciones sin dinamismo ni interacción significativa. Esta presentación pasiva del contenido suele generar una experiencia superficial, donde el turista actúa como receptor y no como participante. 

La interactividad, entendida como la capacidad de respuesta y personalización en función de las decisiones del usuario, es fundamental para mantener su atención y fomentar una exploración más profunda. Diversos estudios demuestran que los elementos de gamificación —como misiones, recompensas, niveles y desafíos— aumentan significativamente el tiempo de uso, el compromiso emocional y la intención de repetir la experiencia turística \cite{Weber2017,Swacha2017}. 

Sin embargo, muchas aplicaciones móviles aún no incorporan estos mecanismos, limitándose a mostrar información sin estimular la curiosidad o la participación activa. En este contexto, surge la necesidad de replantear el diseño de las apps turísticas, integrando dinámicas lúdicas que conviertan al visitante en protagonista de su recorrido.

\subsection{Retención de conocimiento y experiencia cultural}
El turismo no solo implica desplazarse físicamente, sino también adquirir y conservar conocimientos sobre la cultura, historia y tradiciones del lugar visitado. Sin embargo, gran parte de la información entregada a los turistas es olvidada rápidamente debido a su presentación poco memorable o a la falta de conexión emocional con el contenido.

El modelo \textit{edutainment}, que combina educación con entretenimiento, ha demostrado ser una estrategia efectiva para mejorar la retención del conocimiento. Aplicaciones que incluyen elementos como narrativas interactivas (storytelling), quizzes, recompensas simbólicas y desafíos temáticos permiten a los usuarios aprender de manera lúdica, estimulando tanto la memoria como la comprensión \cite{Frontiers2024,Architecture2019}. 

Además, al fomentar la exploración activa, se fortalece la apropiación del patrimonio cultural, promoviendo un turismo más consciente y enriquecedor. En contraste, las soluciones tradicionales tienden a reproducir esquemas pasivos de aprendizaje que no se adaptan a las nuevas generaciones de viajeros digitales.

\subsection{Ausencia de análisis espacial para rutas optimizadas}
La planificación de recorridos turísticos es una de las principales funciones que los visitantes esperan de una aplicación de turismo. Sin embargo, muchas soluciones actuales se limitan a ofrecer mapas básicos o listas de lugares sin una lógica de optimización espacial. Esta carencia se traduce en recorridos poco eficientes, repetitivos o desordenados, que desperdician tiempo y recursos del turista.

La incorporación de sistemas de información geográfica (GIS) permite diseñar rutas inteligentes, considerando variables como distancias, tiempos de visita, preferencias personales, condiciones climáticas o accesibilidad. Además, el enfoque de \textit{participatory sensing}, donde los usuarios contribuyen con información en tiempo real sobre el estado de los sitios, congestión o eventos locales, incrementa aún más la pertinencia y adaptabilidad de las rutas sugeridas \cite{Arxiv2025,Tandfonline2025}.

El turismo moderno requiere soluciones que integren estos análisis espaciales dinámicos para ofrecer experiencias más fluidas, seguras y personalizadas, especialmente en destinos con múltiples puntos de interés dispersos o poco señalizados.

\subsection{Brillo visual limitado y jerarquización subjetiva de contenido}
Uno de los desafíos en la interfaz de usuario de muchas aplicaciones turísticas es la presentación homogénea de los puntos de interés. Al no establecer jerarquías visuales claras, el usuario puede sentirse desorientado o abrumado por la cantidad de información, sin una guía clara sobre qué lugares priorizar o explorar primero.

Investigaciones en diseño de interfaces gamificadas proponen el uso de elementos visuales diferenciadores, como colores intensos, íconos distintivos, efectos de brillo o emblemas de logro (badges), para destacar sitios relevantes en función de su valor cultural, calificaciones de otros usuarios o logros personales del turista \cite{Swacha2017,TravelPlot2018}. 

Estas señales visuales no solo mejoran la usabilidad, sino que también incentivan comportamientos deseados, como visitar lugares menos conocidos pero altamente valorados o completar rutas culturales completas. En contraste, las plataformas sin esta jerarquización conducen a experiencias menos estructuradas y menos significativas.

\subsection{Limitaciones locales}
En muchas regiones con alto potencial turístico, especialmente en áreas rurales o en vías de desarrollo, la infraestructura tecnológica es limitada o inexistente. Esto incluye desde la falta de conectividad móvil adecuada hasta la ausencia total de aplicaciones que orienten al visitante. En estos contextos, los turistas dependen exclusivamente de guías humanos o de la información informal obtenida de boca en boca.

Esta situación genera desigualdades en el acceso al conocimiento, ya que no todos los visitantes reciben la misma calidad o cantidad de información. Asimismo, la dependencia de recursos humanos restringe la escalabilidad del turismo y puede sobrecargar a los actores locales, afectando la sostenibilidad de la actividad \cite{ToARist2020,Intech2019,TourismMobility2021}.

El desarrollo de soluciones digitales adaptadas a estas realidades —por ejemplo, aplicaciones que funcionen offline, que utilicen contenido precargado o que se apoyen en redes comunitarias— es clave para cerrar la brecha tecnológica y fomentar un turismo más inclusivo, autónomo y enriquecedor para todos los perfiles de visitantes.

\newpage
\renewcommand{\refname}{Referencias}
\begin{thebibliography}{99}
\bibitem{Xu2015} Xu, F., Buhalis, D., \& Weber, J. (2015). Gamification in tourism. In \textit{Information and Communication Technologies in Tourism 2015} (pp. 525–537). Springer.
\bibitem{Pradhan2023} Pradhan, P. M., \& Ghimire, M. (2023). Mobile applications for sustainable tourism: a review. \textit{Sustainability}, 15(1), 12.
\bibitem{Weber2017} Weber, J., \& Rech, J. (2017). \textit{Learning and enjoyment in location-based augmented reality applications}. Springer.
\bibitem{Emerald2020} Hernández-Muñoz, M. (2020). Interactive storytelling for tourism heritage: a systematic review. \textit{Journal of Heritage Tourism}, 15(2), 168–183.
\bibitem{Fishoeder2018} Fischoeder, M. (2018). Tourist behavior and information retention in gamified mobile experiences. \textit{Tourism Management}, 65, 72–85.
\bibitem{Architecture2019} Chung, N., Lee, H., \& Han, H. (2019). Heritage tourism in the mobile era: the role of gamification. \textit{Journal of Architecture and Planning}, 84(6), 27–34.
\bibitem{Frontiers2024} Martínez, A., \& Sánchez, L. (2024). Edutainment in mobile heritage tourism: a user-centered approach. \textit{Frontiers in Psychology}, 15, 1883.
\bibitem{Swacha2017} Swacha, J. (2017). Design of a gamification framework for tourist mobile applications. \textit{Procedia Computer Science}, 104, 550–555.
\bibitem{Xu2014} Xu, F., \& Buhalis, D. (2014). Big data analytics in tourism. \textit{Tourism Review}, 69(3), 1–4.
\bibitem{ToARist2020} Ali, A., \& Frew, A. (2020). Absence of digital tools in tourism support in rural areas. \textit{Tourism and Rural Development}, 2(1), 91–105.
\bibitem{Intech2019} Rusu, D., \& Orboi, M. D. (2019). ICT for cultural tourism development in less developed regions. \textit{IntechOpen}.
\bibitem{Arxiv2025} Chen, H., \& Zhao, Y. (2025). Optimizing Tourist Routes Using GIS and AI Models. \textit{arXiv preprint arXiv:2501.01001}.
\bibitem{Tandfonline2025} Liu, M., et al. (2025). Participatory GIS for tourism planning. \textit{Journal of Spatial Science}. https://doi.org/10.1080/14498596.2025.101002
\bibitem{Kawanaka2021} Kawanaka, H. (2021). Visualizing tourist preferences with GIS-based interfaces. \textit{Information Systems in Tourism}, 43(3), 243–257.
\bibitem{TravelPlot2018} Kim, J., \& Kim, H. (2018). TravelPlot: Location-based gamification for tourism. \textit{Entertainment Computing}, 25, 44–52.
\bibitem{TourismMobility2021} García, S., \& Yáñez, J. (2021). Barriers to digital mobility in tourism: the case of isolated regions. \textit{Tourism, Technology and Mobility}, 2(3), 35–49.
\bibitem{NatureGIS2022} Rojas, J., \& Becker, A. (2022). GIS-based tourism analysis for protected areas. \textit{Nature Sustainability}, 5(4), 389–397.
\end{thebibliography}

\end{document}
